\Needspace*{20\baselineskip}
\section{The operator and its normalisation}\label{sec:operator}

This section fixes the operator once, from its definition by Ashtekar and
Lewandowski, in the normalisation in which Brunnemann and Rideout published their
numbers, and checks the normalisation against those numbers.

\subsection{From the definition to a matrix}

Throughout, $J_I=(J_I^1,J_I^2,J_I^3)$ is the Hermitian angular momentum operator of
edge $I$, with $[J_I^a,J_I^b]=i\epsilon^{abc}J_I^c$, acting in the spin-$j_I$
representation; operators on different edges commute. Every edge carries a positive
spin, $j_I\ge\frac12$, and all statements below assume this: an edge of spin $0$ carries
$J_I=0$, so it is not an edge of the vertex, and a five-valent tuple with a spin-$0$ entry
is really a vertex of lower valence. (Allowing such entries would change the constants:
at $t=1$ the tuple $(\frac12;\frac12,\frac12;0,\frac12)$ is the four-valent vertex with
four spins $\frac12$ and has $\qm=2\sqrt3$, below the equal-weight constant of
\cref{sec:tone}.) A \emph{half-integer} is an odd multiple of $\frac12$; when we contrast
integer and half-integer spins, ``half-integer'' excludes the integers. The volume of a region acts
on a spin network through its vertices. At a vertex with edges $e_1,\dots,e_N$,
Ashtekar and Lewandowski~\cite[eq.~(4.13)]{AL1998} define
\[
  \hat V=\kappa_0\sqrt{|\hat q_v|},\qquad
  \hat q_v=\frac1{48}\,\epsilon_{ijk}\sum_{e,e',e''}\epsilon(e,e',e'')\,J^i_eJ^j_{e'}J^k_{e''},
\]
with the sum over ordered triples of edges at the vertex, $\epsilon(e,e',e'')$ the
sign of the determinant of their tangent vectors ($0$ if they are linearly
dependent), and $\kappa_0$ an undetermined constant. Brunnemann and
Rideout~\cite[eqs.~(2.24)--(2.26)]{BR2008b} write the same operator with the
building block
\begin{equation}\label{eq:qhat}
  \hat q_{IJK}:=\bigl[(J_I+J_J)^2,(J_J+J_K)^2\bigr]=-4i\,J_I\cdot(J_J\times J_K)
\end{equation}
and the matrix $Q:=\sum_{I<J<K\le N}\epsilon(IJK)\,\hat q_{IJK}$. Since the summand of
$\hat q_v$ is symmetric under permutations of the three edges,
\[
  \hat q_v=\frac18\sum_{I<J<K}\epsilon(IJK)\,J_I\cdot(J_J\times J_K)
  =\frac18\sum_{I<J<K}\epsilon(IJK)\,\frac i4\,\hat q_{IJK}=\frac{i}{32}\,Q,
\]
and $\hat V=\ell^3\sqrt{|Z|\,|Q|}$ with all constants (including $\kappa_0$, the
Immirzi parameter and $\frac1{32}$) collected in $Z$.
Brunnemann and Rideout set $Z=1$ and publish $\lambda_V:=\sqrt{|\lambda_Q|}$, where
$\lambda_Q$ runs over the eigenvalues of $Q$.

\paragraph{Normalisation used in this paper.}
We write $q$ for an eigenvalue modulus of $Q$ (equivalently of the Hermitian matrix
$iQ$) and $\qm$ for the smallest nonzero one. So $\lambda_V=\sqrt q$, and the
volume eigenvalue is $\ell^3\sqrt{|Z|q}$. A volume gap means $\inf\qm>0$, the
infimum taken, with the embedding fixed, over the assignments of positive spins for which
the operator is not zero on the invariant space.

\subsection{Gauge invariance and the reduced signs}

On gauge-invariant states $\sum_IJ_I=0$, so the last edge can be eliminated:
$J_N=-\sum_{L<N}J_L$ acting on $\Inv(j_1\otimes\cdots\otimes j_N)$, and since $J_N$
commutes with the other edges, $J_I\cdot(J_J\times J_N)=-\sum_{L\ne I,J,N}J_I\cdot
(J_J\times J_L)$. Collecting terms gives Brunnemann and Rideout's
formula~\cite[eq.~(4.12)]{BR2008b},
\begin{equation}\label{eq:sigma}
  Q=\sum_{I<J<K<N}\sigma(IJK)\,\hat q_{IJK},\qquad
  \sigma(IJK)=\epsilon(IJK)-\epsilon(IJN)+\epsilon(IKN)-\epsilon(JKN).
\end{equation}
When no three tangents are coplanar, each $\sigma$ lies in $\{0,\pm2,\pm4\}$. At valence five we list
$\sigma=(\sigma_{123},\sigma_{124},\sigma_{134},\sigma_{234})$ in their order.
In a real recoupling basis each $\hat q_{IJK}$ is a real antisymmetric matrix, so the
spectrum of $iQ$ is symmetric under $q\mapsto-q$ and only moduli matter.

\paragraph{Labelled and sorted spins.}
An embedding fixes $\epsilon$, and we may assign any positive spins to its edges; this is the
\emph{labelled} convention, and statements ``for all spins'' are statements over all
assignments. It is the convention in which a statement about an embedded vertex is a
statement about the spin-network states on that one embedded graph, whatever spins
they carry. Brunnemann and Rideout's 2007 papers instead sort the spins,
$j_1\le\dots\le j_5$, label the edges accordingly, and attach each $\sigma$ to the
sorted tuple~\cite{BR2008b}. A relabelling that sorts the spins permutes the signs,
so one embedding's spectrum is spread over several sorted $\sigma$; their 2010
paper regroups the data by relabelling classes~\cite{BR2010}. The infimum over sorted
spins with $\sigma$ fixed is never smaller than the labelled infimum of the same
embedding, and we note below the two places where they differ.

\subsection{The four-valent matrix}

At valence four only $\sigma(123)\in\{0,\pm2,\pm4\}$ survives and
$Q=-4i\sigma(123)\,J_1\cdot(J_2\times J_3)$. In the basis $|k\rangle$ in which edges
$1$ and $2$ couple to spin $k$, the operator $J_1\cdot(J_2\times J_3)$ is a
zero-diagonal Jacobi matrix with $|\langle k-1|\,J_1\cdot(J_2\times J_3)\,|k\rangle|=a_k$,
\begin{equation}\label{eq:ak}
  a_k^2=\frac{\Delta(j_1,j_2;k)\,\Delta(j_3,j_4;k)}{16\,(4k^2-1)},\qquad
  \Delta(a,b;k)=\bigl((a+b+1)^2-k^2\bigr)\bigl(k^2-(a-b)^2\bigr),
\end{equation}
for $\max(|j_1-j_2|,|j_3-j_4|)<k\le\min(j_1+j_2,j_3+j_4)$. The operator
$J_1\cdot(J_2\times J_3)$ was proposed by Chakrabarti~\cite{Chakrabarti1964} as an operator
whose eigenvalues label the states of three coupled angular momenta, noting that its classical
analogue is the volume generated by the three angular momentum vectors~\cite[p.~302]{Chakrabarti1964}; he wrote it as $\frac i4$
times a commutator of two pair Casimirs, as in \eqref{eq:qhat}~\cite[eq.~(2.4)]{Chakrabarti1964},
and, following a remark of Racah, gave its matrix elements in the coupled basis as a
product of a $6j$ and a $9j$ symbol~\cite[p.~327]{Chakrabarti1964}. L\'evy-Leblond and
L\'evy-Nahas~\cite{LLLN1965} studied it as the generator of a symmetric coupling scheme.
The closed product form \eqref{eq:ak} of the matrix elements is in Brunnemann and
Thiemann~\cite[eq.~(67)]{BT2006}, and, in the form of
a product of two Heron areas
($\Delta=16\,\mathrm{area}^2$ of the triangle with sides $a+\frac12,b+\frac12,k$),
Aquilanti, Marinelli and Marzuoli~\cite[eq.~(2)]{AMM2013}. We recomputed it twice:
by brute force from the generators on the invariant subspace of the tensor product,
and from Racah's formula for the $6j$ symbols in exact arithmetic
(\cref{sec:computer}).

The same four-valent blocks reappear at every valence, because a triple product
$J_I\cdot(J_J\times J_K)$ commutes with the total spin of the three edges and with
every coupling of the remaining ones; on each block it is the four-valent matrix
\eqref{eq:ak} with the fourth spin replaced by the total.

\subsection{Reproduction of published numbers}

With our own code (two independent constructions, \cref{sec:computer}) we
reproduce the numbers that pin down the normalisation, listed in the table below. For
Bianchi and Haggard's column the eigenvalues are those of $J_1\cdot(J_2\times J_3)$,
written $\lambda$ as in \cref{sec:fourvalent}; elsewhere they are the $q$ just
defined.

\begin{center}\small
\begin{tabular}{>{\raggedright\arraybackslash}p{2.9cm}>{\raggedright\arraybackslash}p{4.8cm}>{\raggedright\arraybackslash}p{7.5cm}}
\toprule
\estabhead{source} & \estabhead{published} & \estabhead{reproduced}\\
\midrule
Bianchi--Haggard~\cite{BH2011}, ``loop gravity'' & $\frac{\sqrt2}3\sqrt{|\lambda|}$, four-valent &
  all printed digits: $0.310$ at $(\frac12,\frac12,\frac12,\frac12)$, $0.396$ at
  $(\frac12,\frac12,1,1)$, $0.464$ at $(\frac12,\frac12,\frac32,\frac32)$, $0.498$ at
  $(\frac12,1,1,\frac32)$, $0$ and $0.620$ at $(1,1,1,1)$, $0.522$ at $(\frac12,\frac12,2,2)$,
  $0.577$ at $(\frac12,1,\frac32,2)$, $0.620$ at $(1,1,1,2)$, and the six values
  $1.828,\dots,6.817$ at $(6,6,6,7)$\\
Brunnemann--Rideout~\cite[\S6.5]{BR2008b} & seven-valent $\lambda_V=0.000597922$, $\sigma$ of their
  Table~10, spins $(1,1,1,1,1,\frac32,\frac52)$ & $\sqrt{\qm}=0.000597922$ ($\qm=3.5751089\times10^{-7}$,
  $\dim\Inv=44$)\\
Brunnemann--Rideout~\cite[Table~5]{BR2008b} & minimising sorted spins of $\sigma=(4,2,2,0)$ at each
  $2j_{\max}$ & all $29$ rows, $2\le2j_{\max}\le30$, by exhaustive sorted search\\
Brunnemann--Rideout~\cite[Table~9]{BR2008b} & six-valent minimisers for
  $\sigma=(2,0,0,0,0,0,0,0,-2,0)$ & every level $2j_{\max}\le9$; e.g.\ $\lambda_V=0.0261984$ at
  $(\frac12,\frac12,\frac52,\frac52,3,3)$, where their fit gives $0.0282$\\
Kisielowski~\cite[Table~2]{Kis2021}, Rovelli--Smolin & e.g.\ $1.201406$ six times at valence five,
  spin one; a zero at spin two & all printed digits, with $V=\sqrt{Q_{RS}/32}$ ($\kappa_0=1$
  in~\cite{AL1998})\\
\bottomrule
\end{tabular}
\end{center}
\noindent Four remarks on the table. At $2j_{\max}=30$ the Table~5 minimum is
already $2.5\times10^{-12}$ of the norm of the matrix, where double precision can no
longer rank the candidates, so we stopped there; some text extractions of the published
PDF misalign the rows of that table from $2j_{\max}=10$ on, so we compare against the page
image. Brunnemann and Rideout later reported a code bug that affected all their
seven-valent results~\cite[\S5]{BR2010}; we use their seven-valent value only as a check
of the normalisation at the stated $\sigma$ and spins, where our own code reproduces it. Brunnemann and
Rideout print no eigenvalues at valence six, only minimising spins (their Table~9 and
Figs.~29--32), and their fit $\lambda_V\approx e^{3.72}e^{-1.215\cdot2j_{\max}}$ cannot be
extrapolated (\cref{sec:v6}). At odd $2j_{\max}$ their six-valent minimisers differ
from the family of \cref{sec:v6}.
